\section{Justificación}\label{sec:just}

\subsection{El gran problema}

Los estudiantes subestiman y reparten mal el tiempo que exige un semestre. Según la norma colombiana, un crédito académico equivale a 48 horas de trabajo en el periodo, es decir, unas 3 horas por semana en un periodo de 16 semanas \parencite{MEN2019}. Con el horario real del estudiante (6 materias, 20 créditos) el trabajo total es de 60 horas semanales; solo 21,8 horas ocurren en clase, así que \textbf{las 38,3 horas restantes son estudio independiente} que la persona debe organizar por su cuenta. Esa cifra no aparece en ningún horario ni calendario, por lo que se suele planear de forma reactiva, según lo que se siente más urgente.

El problema se agrava cuando coinciden entregas. En el mismo horario, el sistema detecta que entre el 17 y el 20 de octubre de 2026 coinciden cuatro entregas (dos talleres, un informe de laboratorio y un parcial difícil). Sin una planeación anticipada, esas entregas se acumulan en los últimos días y compiten por el mismo tiempo libre.

\subsection{Por qué es importante resolverlo}

\begin{itemize}
    \item \textbf{Bienestar.} Concentrar el trabajo en pocos días produce jornadas excesivas y estrés. Un plan que reparte la carga y respeta el descanso es más sostenible.
    \item \textbf{Calidad del aprendizaje y de las entregas.} Trabajar con anticipación deja margen para corregir y estudiar, en lugar de entregar lo que se alcanzó a hacer.
    \item \textbf{Visibilidad de la carga real.} Convertir créditos en horas hace explícito algo que hoy permanece implícito y permite decidir con información, por ejemplo al pedir una prórroga a tiempo.
    \item \textbf{Alcance.} Cualquier estudiante tiene el mismo problema, y la herramienta funciona con los datos que ya tiene (su horario), sin requerir información de la institución.
\end{itemize}

El sistema no pretende reemplazar la decisión de la persona: recomienda, explica y deja que el estudiante ajuste.
